\uuid{Vr0o}
\exo7id{4846}
\titre{exo7 4846}
\auteur{quercia}
\organisation{exo7}
\datecreate{2010-03-16}
\isIndication{false}
\isCorrection{false}
\chapitre{Topologie}
\sousChapitre{Espaces complets}
\cpge{mp-mpi.topologie-espaces-normes.a, pc.espaces-vectoriels-normes.a, psi.espaces-vectoriels-normes.a, mp-mpi.topologie-espaces-normes.b, pc.espaces-vectoriels-normes.b, psi.espaces-vectoriels-normes.b, mp-mpi.topologie-espaces-normes.e, pc.espaces-vectoriels-normes.e, psi.espaces-vectoriels-normes.e}
\module{Analyse}
\niveau{L2}
\difficulte{}

\contenu{
\question{
Soit $E = \{ \text{fonctions lipchitziennes } f:\R \to \R \}$.
Pour $f \in E$, on pose
$\|f\| = |f(0)| + \sup\limits_{x\ne y} \left|\frac{f(x)-f(y)}{x-y}\right|$.

Montrer que $E$ est complet.
}
}
